\section{The original two-test connection at the certified determinant zero}
\label{sec:schur-connection-support}

\paragraph{SC1. Original objects and the certified point.}
We keep the original Rodgers--Tao heat family and clock:
\[
 H_t(Z)=\int_0^\infty e^{tu^2}\Phi(u)\cos(Zu)\,du,\qquad
 \partial_tH_t=-\partial_Z^2H_t,\qquad
 H_0(Z)=\frac18\xi_R\!\left(\frac12+\frac{iZ}{2}\right).
 \tag{SC1}
\]
The author source is Brad Rodgers and Terence Tao,
\href{https://arxiv.org/abs/1801.05914v5}{\emph{The de Bruijn--Newman
constant is non-negative}}, version 5, author equations
\texttt{phidef}, \texttt{htdef}, \texttt{hoz}, and \texttt{sas}.
The retained original TeX has canonical source ID
\texttt{RODGERS-TAO-ARXIV-1801.05914-v5}.
Only those defining passages are used here. The complete-root
receiving proof is WR1--5; the fixed tests and support carrier are
NC1--2 and NC21, retained with this component. In particular,
\[
 \begin{gathered}
 f_*(Z)=\sum_{i=0}^{18}c_iZ^{-2i-1},\quad
 c_i=1000^i n_i/10^{120},\quad g_0(Z)=Z^{-1},\\
 Q=\sum_{i,j=0}^{18}c_ic_jS_{i+j+1},\qquad
 B=\sum_{i=0}^{18}c_iS_{i+1},\qquad A=S_1=M_1/M_0>0 .
 \end{gathered}\tag{SC2}
\]
The integers \(n_i\) are the unrounded supplied vector, retained in
\texttt{rational\_witness.json} with SHA256
\path{85314e9fb3115647b7f03a5e00f1b5512c4b23c9280a753f3cf2cababdfed4ac}.
All traces here receive the entire original zero divisor. We use
the ordered original basis \((f_*,g_0)\) of \(V\); its independence
also follows from \(c_{18}=10^{54}\ne0\).

Write
\[
 G(t)=\begin{pmatrix}Q&B\\B&A\end{pmatrix},\quad
 \Delta(t)=\det G=QA-B^2,\quad
 \beta=B/A,\quad S=Q-B^2/A,\quad \Delta=AS.
 \tag{SC3}
\]
The certified calculation in
\path{CERTIFIED_SCHUR_TRANSITION.json} and
\path{certify_schur_transition.py} proves that \(\Delta\) has
one simple zero \(t_d=-2+d_d\) on \([-2,-2+10^{-20}]\), with
\[
 \frac{862686317576371313926327932467260436}{2^{120}10^{20}}
 <d_d<
 \frac{862686317576371313926327932467260437}{2^{120}10^{20}}.
 \tag{SC4}
\]
The proof uses the complete theta interval bounds from the
companion calculation: \(\Delta''>0\) on the interval,
\(\Delta'(-2)>0\), strict opposite endpoint signs, and 120
rational bisections with strict signs. This component receives
that established point; it does not repeat the quadrature.
The same receipt proves \(A_d>0\), \(B_d>0\), and
\[
 \begin{split}
 7.34673724543500384324\,10^{-26}
 &<s_1:=S'(t_d)=\Delta'(t_d)/A_d
 <7.34673724543500384326\,10^{-26},\\
 7.51503947413203068735\,10^{-23}
 &<\beta_d<7.51503947413203068736\,10^{-23}.
 \end{split}\tag{SC5}
\]
These bounds are outward enlargements of the retained exact binary
endpoints. Subscripts \(d\) below denote evaluation at this point.

The local coordinate is the explicit translation
\(\zeta=t-t_d\), so \(d\zeta=dt\); the physical time and the
arithmetic coordinate \(s=\frac12+iZ/2\) are unchanged.
NC1--4 give real analytic entries and their holomorphic extensions
near \(t_d\). Shrink a complex disk so that \(A\ne0\) throughout
and \(S\) has only its simple zero at \(\zeta=0\).
For complex \(\zeta\) we use the holomorphic symmetric bilinear
extension \(v^TGw\). On the real interval, \(G\) is real symmetric
and defines the original Hermitian form \(v^*Gw\). These two uses
of transpose and conjugate transpose will be kept explicit.

\paragraph{SC2. The exact Schur morphism, including its derivative.}
Define
\[
 P=\begin{pmatrix}1&0\\-\beta&1\end{pmatrix},\qquad
 D=\begin{pmatrix}S&0\\0&A\end{pmatrix},\qquad
 C=P^{-1}P'=\begin{pmatrix}0&0\\-\beta'&0\end{pmatrix}.
 \tag{SC6}
\]
Direct multiplication gives \(P^TGP=D\) and \(\det P=1\).
Thus \(P\) is a holomorphic isometry from the Schur presentation
\((\mathbb C^2,D)\) to the original presentation
\((V,G)\), including the degenerate central fibre. The first
Schur basis vector maps to \(f_*-\beta g_0\); the second maps to
\(g_0\). No scalar factor has been removed from the metric.

On the punctured disk, set
\[
 \Omega=\Gamma\,dt,\qquad
 \Gamma=\frac12G^{-1}G'
 =\frac1{2AS}
 \begin{pmatrix}
 A Q'-B B'&A B'-B A'\\
 Q B'-B Q'&Q A'-B B'
 \end{pmatrix}.
 \tag{SC7}
\]
The horizontal equation is \(X'=-\Gamma X\).
Because \(G^T=G\) and \((G')^T=G'\),
\[
 G'=\Gamma^TG+G\Gamma.
 \tag{SC8}
\]
For horizontal column solutions \(x,y\), differentiating
\(x^TGy\) therefore gives zero. For real time and arbitrary complex
column solutions the same argument, with \(\Gamma^*=\Gamma^T\),
proves preservation of \(x^*Gy\). The connection is flat on the
punctured complex disk: both \(d\Omega\) and
\(\Omega\wedge\Omega\) vanish, since every factor is a multiple
of \(dt\).

For the change of variables \(x=Py\), the full transformed
coefficient is
\[
 \begin{split}
 \widetilde\Gamma
 &=P^{-1}\Gamma P+P^{-1}P'\\
 &=\frac12D^{-1}D'+\frac12(C-D^{-1}C^TD)\\
 &=\begin{pmatrix}
       S'/(2S)&A\beta'/(2S)\\
       -\beta'/2&A'/(2A)
    \end{pmatrix}.
 \end{split}\tag{SC9}
\]
To prove the middle equality, differentiate \(D=P^TGP\):
\[
 P^TG'P=D'-C^TD-DC,\qquad
 P^{-1}G^{-1}G'P=D^{-1}(D'-C^TD-DC).
\]
Add \(C\) after division by two. The correction
\(K=(C-D^{-1}C^TD)/2\) satisfies \(K^TD+DK=0\).
Consequently both displayed terms contribute to the same exact
compatibility equation \(D'=\widetilde\Gamma^TD+
D\widetilde\Gamma\). Deleting \(K\) changes the connection; it
is not the gauge transformation of SC7 unless \(\beta'=0\).

\paragraph{SC3. The residue in the original basis.}
Put
\[
 r=\binom{1}{-\beta_d},\qquad
 \ell=\binom{\beta_d}{1},\qquad
 G_d=A_d\ell\ell^T,\qquad
 \operatorname{adj}G_d=A_drr^T .
 \tag{SC10}
\]
The determinant derivative is \(\Delta'_d=A_ds_1\).
Dividing the adjugate formula for \(G^{-1}\) by its simple
determinant zero gives
\[
 R:=\operatorname{Res}_{t=t_d}\Omega
   =\frac{r r^TG'_d}{2s_1}.
 \tag{SC11}
\]
One also has \(r^TG'_dr=s_1\). Indeed
\((1,-\beta)G(1,-\beta)^T=S\), and the two differentiated-vector
terms vanish at \(t_d\), because \(G_dr=0\). It follows that
\[
 R^2=\tfrac12R,\quad
 \operatorname{tr}R=\tfrac12,\quad
 \operatorname{im}R=\mathbb Cr=\ker G_d,\quad
 \ker R=\ker(r^TG'_d).
 \tag{SC12}
\]
Every assertion follows directly from SC11 and
\(r^TG'_dr=s_1\ne0\): \(R\) has rank one, \(Rr=r/2\),
and its defining covector is nonzero.

For an explicit second line, set
\[
 a=\frac{A_d\beta'_d}{s_1},\qquad
 v=\binom{-a}{1+\beta_da},\qquad
 T=(r\ \ v)=
 \begin{pmatrix}1&-a\\-\beta_d&1+\beta_da\end{pmatrix}.
 \tag{SC13}
\]
The identity
\[
 \frac{r^TG'_d}{s_1}=(1+\beta_da,\ a)
 \tag{SC14}
\]
follows by differentiating \(B=A\beta\) and
\(Q=S+A\beta^2\). Hence \(Rv=0\), \(\det T=1\), and
\[
 T^{-1}RT=\Lambda=\begin{pmatrix}1/2&0\\0&0\end{pmatrix},
 \qquad v^TG_dv=A_d.
 \tag{SC15}
\]
In the Schur presentation,
\[
 \widetilde R=
 \begin{pmatrix}1/2&a/2\\0&0\end{pmatrix},\qquad
 R=P_d\widetilde RP_d^{-1}
 =\frac12\begin{pmatrix}
 1+\beta_da&a\\-\beta_d(1+\beta_da)&-\beta_da
 \end{pmatrix}.
 \tag{SC16}
\]
The derivative \(P^{-1}P'\) is analytic, so it contributes no
residue, even though it is essential to the full connection.
For the actual input, the certified residue is enclosed by
\[
\begin{aligned}
 R_{11}&\in(0.56568826983857144053,\,0.56568826983857144055),\\
 R_{12}&\in(8.74090815686078356647,\,8.74090815686078356648)10^{20},\\
 R_{21}&\in(-4.25116967789031619450,\,-4.25116967789031619449)10^{-23},\\
 R_{22}&\in(-0.06568826983857144054,\,-0.06568826983857144053).
\end{aligned}
 \tag{SC17}
\]
The exact rank, trace, eigenvalues, and lines are SC11--16;
independent rounding of matrix entries is not used to prove them.

\paragraph{SC4. A convergent gauge retaining every analytic coefficient.}
SC7 and the simple zero imply
\(\Gamma=R/\zeta+J(\zeta)\), with \(J\) holomorphic at zero.
Write \(J=\sum_{j\ge0}J_j\zeta^j\).
There is a unique holomorphic matrix \(F\), with \(F(0)=I\),
satisfying
\[
 F^{-1}\Gamma F+F^{-1}F'=\frac R\zeta .
 \tag{SC18}
\]
Here and below the disk may be made smaller without altering the
physical parameter or the matrix entries.

We prove existence, convergence, and invertibility. Substitution
of \(F=\sum_{k\ge0}F_k\zeta^k\), \(F_0=I\), into
\(\zeta F'=FR-RF-\zeta JF\) gives the exact recurrence
\[
 (kI+\operatorname{ad}_R)F_k
   =-\sum_{j=0}^{k-1}J_jF_{k-1-j},\qquad k\ge1,
 \quad \operatorname{ad}_R(U)=RU-UR .
 \tag{SC19}
\]
In the \(T\) basis, each entry on the left is multiplied by
\(k+\lambda_i-\lambda_j\), where
\((\lambda_1,\lambda_2)=(1/2,0)\). The four multipliers are
\(k,k+1/2,k-1/2,k\), none zero. This proves formal existence
and uniqueness.

For convergence use
\(\|U\|_*=\|T^{-1}UT\|_1\), where \(\|\cdot\|_1\) is the
induced column-sum matrix norm. This norm is submultiplicative,
and \(\|I\|_*=1\). Entrywise division by the four multipliers
shows
\[
 \|(kI+\operatorname{ad}_R)^{-1}U\|_*
 \le\frac{\|U\|_*}{k-1/2}\le\frac2k\|U\|_*.
 \tag{SC20}
\]
Choose \(\rho>0\) such that \(J\) is holomorphic on a neighborhood
of the closed disk of radius \(\rho\), and let
\(M=\max_{|\zeta|=\rho}\|J(\zeta)\|_*\).
Cauchy's coefficient formula gives
\(\|J_j\|_*\le M\rho^{-j}\). Therefore
\[
 \|F_k\|_*\le\frac2k\sum_{j=0}^{k-1}
 M\rho^{-j}\|F_{k-1-j}\|_*.
\]
The nonnegative majorant coefficients \(g_k\) defined by equality
and \(g_0=1\) are the Taylor coefficients of
\[
 g(z)=\exp\!\left(2M\int_0^z\frac{dw}{1-w/\rho}\right)
     =(1-z/\rho)^{-2M\rho}.
 \tag{SC21}
\]
This function is analytic for \(|z|<\rho\), with nonnegative
Taylor coefficients, including the case \(M=0\).
Induction gives \(\|F_k\|_*\le g_k\); hence \(F\) converges
absolutely and locally uniformly there. Its differentiated
series satisfies SC18 wherever it is invertible, initially on
a disk because \(F(0)=I\). Taking the trace of that equation yields
\[
 (\det F)'=-(\operatorname{tr}J)\det F,\qquad
 \det F(\zeta)=
 \exp\!\left(-\int_0^\zeta\operatorname{tr}J(w)\,dw\right).
 \tag{SC22}
\]
The equality first holds on that smaller disk and then throughout
the disk by the identity theorem, since both sides are analytic.
Its right side never vanishes. Thus \(F\) is invertible wherever
the preceding convergence construction is used. All coefficients
of \(J\) enter SC19; no analytic remainder has been dropped.

\paragraph{SC5. An exact metric identity and the original real inertia.}
Set \(E=FT\). Then SC18 and SC15 give
\[
 E^{-1}\Gamma E+E^{-1}E'=\frac{\Lambda}{\zeta}.
 \tag{SC23}
\]
The following identity retains both original constants:
\[
 E(\zeta)^T G(t_d+\zeta)E(\zeta)
   =\begin{pmatrix}s_1\zeta&0\\0&A_d\end{pmatrix}.
 \tag{SC24}
\]
To prove it, denote the left side by \(K\).
It is holomorphic at zero. Differentiate using SC8 and SC23:
\(\zeta K'=\Lambda^TK+K\Lambda\).
If \(K_{ij}=\sum_{n\ge0}k_{ij,n}\zeta^n\), then
\[
 (n-\lambda_i-\lambda_j)k_{ij,n}=0.
\]
Thus \(K_{11}\) is a constant multiple of \(\zeta\),
\(K_{22}\) is constant, and \(K_{12}=K_{21}=0\), because
\(1/2\) is not a nonnegative integer. At zero,
\(E(0)=T\); hence \(K_{22}(0)=v^TG_dv=A_d\).
For the first entry, the differentiated-vector terms vanish
because \(G_dr=0\), giving \(K'_{11}(0)=r^TG'_dr=s_1\).
This proves SC24.

Reality of the input and SC19 shows that every \(F_k\) is real.
For real \(\zeta\), SC24 is therefore also an exact Hermitian
congruence, with \(E^*=E^T\). Since \(s_1,A_d>0\), the
original two-test form has one negative and one positive
eigenvalue for sufficiently small negative \(\zeta\), one
zero and one positive eigenvalue at \(\zeta=0\), and two
positive eigenvalues for sufficiently small positive \(\zeta\).
These are statements about the fixed subspace
\(\operatorname{span}_{\mathbb C}(f_*,g_0)\). No positivity
statement for the full nineteen-dimensional test space follows
from this two-dimensional congruence.

The complete morphism from the diagonal presentation in SC24
to the Schur presentation is \(Y=P^{-1}E\). Direct substitution
of SC9 and SC23 proves
\[
 Y^TDY=\operatorname{diag}(s_1\zeta,A_d),\qquad
 Y^{-1}\widetilde\Gamma Y+Y^{-1}Y'=\Lambda/\zeta.
 \tag{SC25}
\]
The composition \(PY=E\) holds on the whole disk, including zero
for the metrics and on the punctured disk for the connections.
This proves both changes of presentation and their exact
composition; it also exhibits where the derivative correction
in SC9 is received.

\paragraph{SC6. Full local monodromy and the based reflection.}
On the logarithmic covering of the punctured disk a horizontal
fundamental matrix is
\[
 X(\zeta)=F(\zeta)\exp(-R\log\zeta).
 \tag{SC26}
\]
Differentiation using SC18 verifies \(X'=-\Gamma X\);
both factors are invertible. A positive meridian adds \(2\pi i\)
to \(\log\zeta\). Since \(R^2=R/2\),
\[
 e^{-2\pi iR}=I-4R,\qquad (I-4R)^2=I.
 \tag{SC27}
\]
For example, expand the exponential on
\(\mathbb C^2=\operatorname{im}R\oplus\ker R\):
it is multiplication by \(e^{-\pi i}=-1\) on the first
line and by \(1\) on the second.

Fix a base point \(\zeta_b\ne0\) in the disk. Parallel transport
around a positive meridian, in the actual original basis at
that point, is
\[
 \mathcal M_b
 =F_b(I-4R)F_b^{-1},\quad
 \operatorname{im}(\mathcal M_b-I)=F_b\mathbb Cr,\quad
 \ker(\mathcal M_b-I)=F_b\mathbb Cv.
 \tag{SC28}
\]
Indeed analytic continuation of SC26 is right multiplication by
\(I-4R\), so the based transport is \(X_b(I-4R)X_b^{-1}\).
The exponential in \(X_b\) commutes with \(R\), yielding SC28.
Thus the original-basis matrix at a nonzero base point involves
the full convergent \(F_b\). Its two invariant lines tend to
\(\mathbb Cr\) and \(\mathbb Cv\) as the base point tends to zero.
Changing the base point by parallel transport conjugates the
matrix, as follows directly from the same fundamental solution.

There is also an exact formula using the original metric.
Put \(u_b=F_br\), \(w_b=F_bv\), and \(G_b=G(t_d+\zeta_b)\).
SC24 gives
\[
 u_b^TG_bu_b=s_1\zeta_b,\qquad
 w_b^TG_bw_b=A_d,\qquad u_b^TG_bw_b=0.
\]
Consequently
\[
 \mathcal M_b
   =I-\frac{2u_bu_b^TG_b}{s_1\zeta_b},\qquad
 \mathcal M_b^TG_b\mathcal M_b=G_b,\qquad
 \det\mathcal M_b=-1 .
 \tag{SC29}
\]
To verify the first equality, apply its right side to the
basis \((u_b,w_b)\): it sends the first vector to its negative
and fixes the second, exactly as SC28. The other identities
follow either by the same basis calculation or by SC8.
For real base points \(u_b,w_b,\mathcal M_b\) are real, so the
orthogonality identity also uses conjugate transpose for the
original Hermitian form.

The fundamental group of the punctured disk is generated by
one meridian, and a loop of winding number \(k\) has holonomy
\(\mathcal M_b^k\); its determinant is \((-1)^k\).
Alternatively \(\operatorname{tr}\Gamma=
\frac12(\log\Delta)'\), and the determinant equation for
horizontal transport gives the same character, since \(\Delta\)
has a simple zero. This also verifies the sign in the
horizontal convention of SC7.

\paragraph{SC7. Supported determinant, radical, and exact quotient.}
Let \(L\) be the original bounded distributive lattice. For a
complex vector space \(W\), retain the full carrier
\[
 M_L(W)=(W\times\{1_L\})\cup(\{0\}\times L).
 \tag{SC30}
\]
Addition is \((w,\lambda)+(z,\mu)=(w+z,\lambda\vee\mu)\);
scalar multiplication by \((c,\nu)\in M_L(\mathbb C)\) is
\((cw,\nu\wedge\lambda)\). The zero is
\(\tau=(0,0_L)\), and \(e=(0,1_L)\) is the top supported
zero. Here \(\tau\) denotes this zero element, and the local
time displacement remains \(\zeta\).
The operations stay in the stated carrier: any amplitude
different from zero forces every required support to be top.
Associativity and the distributive laws follow from those of
vector addition, scalar multiplication, and the distributivity
of \(L\). Complex conjugation acts on the amplitude and fixes
the support.

For real \(t\), the complete sesquilinear lift of the form is
\[
 \widehat Q_t((z,\lambda),(w,\mu))
  =(z^*G(t)w,\lambda\wedge\mu).
 \tag{SC31}
\]
Additivity uses
\((\lambda\vee\nu)\wedge\mu
=(\lambda\wedge\mu)\vee(\nu\wedge\mu)\);
scalar compatibility uses associativity of meet and the
original sesquilinearity. Thus SC31 proves the lift, with the
amplitude projection commuting with the original form.
For complex time the same formula with \(z^T\) supplies the
holomorphic bilinear version.

The entries of the original Gram array have top support.
Taking its determinant in that top fibre gives
\[
 \widehat\Delta(t)
 =(Q,1_L)(A,1_L)+(-1,1_L)(B,1_L)^2
 =(\Delta(t),1_L).
 \tag{SC32}
\]
Hence \(\widehat\Delta(t_d)=e\). If \(0_L\ne1_L\), this is
distinct from \(\tau\); if \(L\) has one element the labels
coincide, exactly as the definition requires. In every case the
amplitude determinant is zero and has no multiplicative
inverse. The supported label does not supply an inverse for
the matrix in SC7 at \(t_d\).

At the central fibre, the exact sequence of ordinary vector
spaces is
\[
 0\longrightarrow\mathbb Cr
 \ \xrightarrow{\ \iota\ }\ V
 \ \xrightarrow{\ \ell^T\ }\ \mathbb C
 \longrightarrow0,\qquad
 G_dz=A_d\ell\,(\ell^Tz).
 \tag{SC33}
\]
Here coordinates in \(V\) refer to the original ordered tests,
and \(G_d\) maps to the coordinate dual space. The map
\(\ell^T\) is surjective, since \(\ell^T(0,1)^T=1\);
its kernel is exactly \(\mathbb Cr\). SC10 proves the
factorization. Thus the image of the Gram map is
\(\mathbb C\ell\), the radical is \(\mathbb Cr\), and the
quotient \(V/\mathbb Cr\) receives the nondegenerate form
\[
 Q_d(z,w)=A_d\,\overline{\ell^Tz}\,(\ell^Tw).
 \tag{SC34}
\]
This is the full quotient form, including \(A_d\).

Every linear map \(U:W\to W'\) has the explicit supported lift
\(\widehat U(w,\lambda)=(Uw,\lambda)\).
It preserves addition and supported scalar multiplication by
the displayed formulas. Composition and amplitude projection
commute with this lift. Therefore SC33, the maps \(P,E,Y\),
and all their inverses where present have exact supported
receivers without losing any zero label.
For the degenerate Gram map the amplitude radical carrier is
\[
 \mathcal N_L=(\mathbb Cr\times\{1_L\})
               \cup(\{0\}\times L).
 \tag{SC35}
\]
The inverse image of the full zero-amplitude carrier under
\(\widehat G_d\) is exactly \(\mathcal N_L\): this follows
from \(\ker G_d=\mathbb Cr\), with each support retained.
For \(c\ne0\), \(\widehat G_d(cr,1_L)=(0,1_L)\), and its
pairing with every top-supported vector is \(e\).
For an argument supported at \(\lambda\), the output support
is the stated meet, including when its amplitude vanishes.
This gives the precise relation between the amplitude radical
and the retained supported zeros. In particular it does not
replace all zero amplitudes by the additive zero \(\tau\).

At each nonzero base point the invertible holonomy
\(\mathcal M_b\) lifts to
\(\widehat{\mathcal M}_b(z,\lambda)=(\mathcal M_bz,\lambda)\).
It preserves SC31 for real base points, because of SC29, and
fixes every labelled zero \((0,\lambda)\). On amplitudes it
reverses \(F_b\mathbb Cr\) and fixes \(F_b\mathbb Cv\).
Its supported determinant is \((-1,1_L)\).

\paragraph{SC8. What the local holonomy receives.}
Equations SC18--29 calculate the complete local connection of
the actual two-test form at the established \(t_d\).
They prove that its local holonomy receives a simple loss of
rank in that fixed Gram space: the residue image is the
central radical, and the nonzero-fibre negative eigenspace of
the holonomy converges to it. The analytic transport of that
line is given by the convergent \(F\), whose coefficients are
SC19, and its exact metric value is \(s_1\zeta\).

The connection is connected to the original heat family by the
explicit complete-root trace map SC2 and its original-clock
derivatives NC3. Its singularity is the zero of the determinant
of those received pairings. This calculation neither
identifies that determinant with \(H_t\) nor asserts a zero of
\(H_t\) at a specified spatial point. The certified physical
time is the strictly negative number in SC4, and remains near
\(-2\) after every displayed change of frame. Thus the local
reflection is established at that two-test transition; no
claim of an RH proof, disproof, or endpoint transition at
\(t=0\) is made.

\paragraph{SC9. Reproducible finite verification.}
The companion \texttt{check\_schur\_connection.py} uses exact
rational operations. It checks SC6--16 on explicit polynomial
analytic metrics with nonconstant \(A,S/\zeta,\beta\); computes
the first eight convergent-gauge coefficients using SC19;
checks the differential equation and SC24 to those orders;
and verifies the reflection identities. These finite checks
test signs, transpose order, noncommuting products, and the
frame derivative correction. SC18--24 supply the separate
all-orders convergence proof. The checker also reads the exact
binary endpoints of the actual receipt to verify every
outward enlargement displayed in SC4--5 and SC17, without
repeating any numerical integration.

\paragraph{SC10. The original nineteen-test inclusion and its full connection.}
The companion complete-theta calculation proves that the
original nineteen-test matrix is nondegenerate throughout
\(I=[-2,-2+10^{-20}]\), with inertia \((18,1,0)\) in the order
(positive, negative, zero). The exact receiver here is
\[
 \begin{gathered}
 \mathcal V=\operatorname{span}_{\mathbb C}
      \{h_i(Z)=Z^{-2i-1}:0\le i\le18\},\qquad
 \mathcal K(t)=(S_{i+j+1}(t))_{0\le i,j\le18},\\
 D_0=\operatorname{diag}(1000^i)_{i=0}^{18},\qquad
 M=D_0\mathcal KD_0,\qquad w_i=n_i/10^{120},\\
 J_Z=(c,e_0),\qquad
 J=D_0^{-1}J_Z=(w,e_0)
 :\mathbb C^2\longrightarrow\mathbb C^{19},\qquad
 J^TMJ=G .
 \end{gathered}\tag{SC36}
\]
These are both original coefficient presentations. The
unscaled coordinate map sends \(z\) to \(\sum z_i h_i\).
The NF1 coordinate map sends \(x\) to
\(\sum x_i1000^i h_i\); the exact morphism between the
coordinates is \(z=D_0x\). Thus \(J_Z\) and \(J\) both give
the original constant inclusion of \((f_*,g_0)\), with no
alteration of spatial variable, heat clock, or rational test.
Since \(w_{18}=1\), \(J\) is injective. We use the NF1
symbols \(M,J\) below. The supplied proof receipt is
\texttt{UNIFORM\_FULL\_FLAG\_RECEIPT.json}: it encloses the
lower eighteen-dimensional fixed-frame perturbation norm
below \(1/20\), and the attained final Schur complement in
\[
 (-2.403366914664958469,\,-2.364761442970871715)\,10^{-46}.
 \tag{SC37}
\]
The companion argument retains all theta, Taylor, and
fifth-derivative remainder bounds. We use its established
nondegeneracy at the actual \(t_d\); these eighteen positive
directions and one negative direction therefore coexist with
the two-test transition proved above.

Because \(\det M(t_d)\ne0\), it remains nonzero on a complex
disk about \(t_d\). The full original-basis connection is
\[
 \Gamma_{\rm full}=\tfrac12M^{-1}M',\qquad
 M'=\Gamma_{\rm full}^TM+M\Gamma_{\rm full}.
 \tag{SC38}
\]
It is holomorphic throughout this disk. Its horizontal
fundamental matrix \(U\), with \(U(0)=I_{19}\), is holomorphic
and invertible there after shrinking the disk. For completeness,
if \(\Gamma_{\rm full}=\sum_{j\ge0}L_j\zeta^j\), its
coefficients are uniquely given by
\[
 nU_n=-\sum_{j=0}^{n-1}L_jU_{n-1-j},\qquad U_0=I_{19}.
 \tag{SC39}
\]
For a submultiplicative matrix norm, Cauchy bounds
\(\|L_j\|\le C\rho^{-j}\) give the convergent majorant
\((1-z/\rho)^{-C\rho}\), by the same coefficient comparison
as SC20--21, now with factor \(1/n\). Thus the series converges.
The determinant satisfies
\((\det U)'=-\operatorname{tr}(\Gamma_{\rm full})\det U\)
and is the exponential of a holomorphic integral. Hence it
does not vanish. This proves that full parallel transport is
\(U(t_2)U(t_1)^{-1}\) and has identity holonomy around every
loop in this disk, including a loop surrounding \(t_d\).
For the unscaled Laurent coordinates the coefficient is
\(\Gamma_{\mathcal K}=\mathcal K^{-1}\mathcal K'/2\).
Since \(D_0\) is constant, differentiation and inversion give
the exact gauge identity
\(\Gamma_{\rm full}=D_0^{-1}\Gamma_{\mathcal K}D_0\).
Thus \(D_0\) intertwines their horizontal equations and
parallel transports, including the identity local holonomy.

\paragraph{SC11. Exact compression and the complementary term.}
On the punctured disk define
\[
 \begin{split}
 J^\dagger&=G^{-1}J^TM:\mathbb C^{19}\longrightarrow\mathbb C^2,\\
 \Pi&=JJ^\dagger:\mathbb C^{19}\longrightarrow\mathbb C^{19},\\
 W&=(I_{19}-\Pi)\Gamma_{\rm full}J .
 \end{split}\tag{SC40}
\]
The dagger denotes this metric adjoint and left inverse; it
is the stated meromorphic holomorphic formula, not a
conjugation of complex time. On real time it is also the
adjoint of the original Hermitian metrics, since \(J\) is
real. From SC36,
\[
 \begin{gathered}
 J^\dagger J=I_2,\qquad \Pi^2=\Pi,\qquad
 \Pi^TM=M\Pi,\\
 \operatorname{im}\Pi=\operatorname{im}J,\qquad
 \ker\Pi=\ker(J^TM)=(\operatorname{im}J)^{\perp_M}.
 \end{gathered}\tag{SC41}
\]
Indeed \(J^\dagger J=G^{-1}G\); idempotence and the image
follow immediately. Since \(G^{-1}\) is symmetric,
\(M\Pi=MJG^{-1}J^TM\) is symmetric. Finally \(J\) is
injective and \(G^{-1}\) invertible, proving the kernel
identity.
In the unscaled presentation define
\(J_Z^\dagger=G^{-1}J_Z^T\mathcal K\) and
\(\Pi_Z=J_ZJ_Z^\dagger\). SC36 gives
\[
 J^\dagger=J_Z^\dagger D_0,\qquad
 \Pi=D_0^{-1}\Pi_ZD_0 .
\]
Thus the projection, its image, kernel, and residue below
are carried to the unscaled original Laurent coordinates
by the same exact constant isomorphism \(D_0\).

Differentiating SC36 uses \(J'=0\). Since
\(M\Gamma_{\rm full}=M'/2\), it gives the exact
compression and decomposition
\[
 \begin{split}
 J^\dagger\Gamma_{\rm full}J
   &=\tfrac12G^{-1}J^TM'J=\tfrac12G^{-1}G'=\Gamma,\\
 \Gamma_{\rm full}J&=J\Gamma+W,\qquad
 J^TMW=0 .
 \end{split}\tag{SC42}
\]
Thus the compressed connection is the orthogonally projected
full connection on this fixed subspace for every nonzero
\(\zeta\). Its unprojected complementary term is the displayed
\(W\); it is not removed from the full equation. For a
compressed horizontal solution \(y'=-\Gamma y\), the
original ambient vector \(x=Jy\) obeys
\[
 x'+\Gamma_{\rm full}x=W y.
 \tag{SC43}
\]
This follows by substitution in SC42. It states precisely
which ambient covariant derivative remains when one uses
compressed horizontal transport.

\paragraph{SC12. Nilpotent projection residue and cancellation.}
Let \(n=Jr\ne0\). From SC10--11 one has the Laurent expansions
\[
 \begin{split}
 G^{-1}&=\frac{rr^T}{s_1\zeta}+\hbox{a holomorphic matrix},\\
 \operatorname{Res}J^\dagger&=\frac{r n^TM_d}{s_1},\\
 K_\Pi:=\operatorname{Res}\Pi&=\frac{n n^TM_d}{s_1}.
 \end{split}\tag{SC44}
\]
In particular
\[
 K_\Pi\ne0,\qquad K_\Pi^2=0,\qquad
 \operatorname{im}K_\Pi=\mathbb Cn,\qquad
 \ker K_\Pi=\ker(n^TM_d).
 \tag{SC45}
\]
To prove these assertions, observe that
\(n^TM_dn=r^TG_dr=0\). Since \(M_d\) is invertible and
\(n\ne0\), the row \(n^TM_d\) is nonzero. The displayed
outer product therefore has rank one and square zero,
with the indicated image and kernel. This also proves that
the pole of \(\Pi\) is a genuine pole, while the ambient
connection remains holomorphic.

The singular parts in SC42 cancel exactly:
\[
 \operatorname{Res}(J\Gamma)=JR,\qquad
 \operatorname{Res}W=-JR,\qquad
 \operatorname{Res}(\Gamma_{\rm full}J)=0 .
 \tag{SC46}
\]
One proof is the exact analytic decomposition SC42.
An independent calculation gives
\[
 K_\Pi(\Gamma_{\rm full})_dJ
  =\frac{n n^TM_d(\Gamma_{\rm full})_dJ}{s_1}
  =\frac{Jr\,r^TJ^TM'_dJ}{2s_1}
  =JR ,
 \tag{SC47}
\]
which proves SC46 directly from the projection's principal
part. Since \(J\) is injective and \(R\ne0\), neither residue
being cancelled is zero.

At the central fibre the exact geometric relation is
\[
 \operatorname{im}J\cap(\operatorname{im}J)^{\perp_{M_d}}
     =\mathbb Cn .
 \tag{SC48}
\]
Indeed \(Jz\) belongs to the intersection exactly when
\(J^TM_dJz=G_dz=0\); SC12 identifies those \(z\) as
\(\mathbb Cr\). The two spaces have dimensions \(2\) and
\(17\), because \(J^TM_d\) has rank two. They have
one-dimensional intersection at \(t_d\), so their sum there
has dimension \(18\). For \(\zeta\ne0\), SC41 gives instead
their direct sum as the full nineteen-dimensional space.
The line \(\mathbb Cn\) is isotropic in the nondegenerate
ambient metric, but is not an ambient radical: \(M_dn\ne0\).
This follows from invertibility, with its exact nonzero
covector already displayed in SC44.
In fact the limiting sum has the precise identification
\[
 \operatorname{im}J+(\operatorname{im}J)^{\perp_{M_d}}
   =\ker(n^TM_d)=\ker K_\Pi .
 \tag{SC49}
\]
For \(Jz\), the row \(n^TM_d\) gives
\(r^TG_dz=0\). It also vanishes on the orthogonal complement
because \(n\in\operatorname{im}J\). Thus the sum lies in that
kernel. Both have dimension eighteen, the latter because
\(n^TM_d\ne0\), proving equality.
Consequently the projection residue induces the exact
isomorphism
\[
 \mathbb C^{19}/\bigl(\operatorname{im}J+
            (\operatorname{im}J)^{\perp_{M_d}}\bigr)
 \longrightarrow\mathbb Cn,\qquad
 [z]\longmapsto \frac{n(n^TM_dz)}{s_1}.
\]
It is well defined and injective by SC49, and surjective
by the image assertion in SC45. This identifies the
one-dimensional quotient determined by the failed
orthogonal direct sum with the central compressed radical.

The projection pole, its nilpotent residue, and the
complementary residue in SC46 are thus the exact maps
connecting the compressed reflection SC28 to the ambient
identity holonomy SC39. The full horizontal equation and
the compressed equation are related by SC43, with a
nonzero complementary meromorphic operator. Removing it
would replace the ambient equation. Every map in SC36 and
SC40--48 has its support-preserving lift as in SC30--35
where its matrix is defined. In particular \(J\) and \(M_d\)
remain defined and injective at the central fibre;
\(\widehat M_d(n,1_L)\) has nonzero amplitude, while its
compression by \(\widehat{J^T}\) has amplitude zero and
top support. This proves the relation of the retained
compressed zero to the nondegenerate original full form.

There is a direct identity between the two parallel transports.
Along any piecewise smooth path in the punctured disk starting
at \(b\), let \(\mathcal P_V(z,b)\) denote the compressed
transport, continued along that path, and put
\(\mathcal P_{\rm full}(z,w)=U(z)U(w)^{-1}\).
The full matrix \(U\) is single valued on the entire disk.
Variation of constants applied to SC43 gives
\[
 \begin{split}
 J\mathcal P_V(z,b)
   ={}&\mathcal P_{\rm full}(z,b)J\\
    &+\int_b^z
      \mathcal P_{\rm full}(z,w)\,
      W(w)\mathcal P_V(w,b)\,dw .
 \end{split}\tag{SC50}
\]
To prove the formula, differentiate
\(U(z)^{-1}J\mathcal P_V(z,b)\). The product rule and SC43
give \(U(z)^{-1}W(z)\mathcal P_V(z,b)\). Integrate on the
chosen path and multiply by \(U(z)\); its value at \(b\)
is \(U(b)^{-1}J\). For a positive meridian \(\gamma\) based
at \(b\), SC50 becomes
\[
 J(\mathcal M_b-I_2)
   =\int_\gamma U(b)U(w)^{-1}W(w)
          \mathcal P_V(w,b)\,dw
   =-4JF_bRF_b^{-1}\ne0 .
 \tag{SC51}
\]
The integrand uses its continued compressed solution along
the path; it is not asserted to be single valued on the
punctured disk. The final matrix is nonzero because \(J\)
is injective, \(F_b\) is invertible, and \(R\ne0\).
This identity identifies the entire difference of the two
holonomies through the retained complementary term.

\begin{figure}[p]
\centering
\begin{tikzpicture}[>=Stealth,
 box/.style={draw,rounded corners,align=center,text width=4.5cm,
 inner sep=8pt}]
\node[box] (diag) at (0,0)
 {Analytic presentation\\
 metric \(\operatorname{diag}(s_1\zeta,A_d)\)\\
 connection \(\Lambda\,dt/\zeta\)};
\node[box] (schur) at (7,0)
 {Schur presentation\\metric \(D=\operatorname{diag}(S,A)\)\\
 connection \(\widetilde\Gamma\,dt\) in SC9};
\node[box] (original) at (3.5,-4)
 {Original tests \((f_*,g_0)\)\\
 metric \(G\)\\ connection \(G^{-1}G'\,dt/2\)};
\draw[->,thick] (diag)--node[above]{\(Y=P^{-1}E\)} (schur);
\draw[->,thick] (diag)--node[left]{\(E=FT\)} (original);
\draw[->,thick] (schur)--node[right]{\(P\)} (original);
\node[align=center,text width=13cm] at (3.5,-6.1)
 {\(PY=E\), \(d\zeta=dt\), and every gauge includes its derivative.\\
 At \(\zeta=0\): metric radical maps to \(\mathbb C(f_*-\beta_dg_0)\).};
\node[box,text width=11.6cm] at (3.5,-8.3)
 {At a punctured-disk base point:
 \(\mathcal M_b=F_b(I-4R)F_b^{-1}\).\\
 \(-1\) line \(F_b\mathbb Cr\), \(+1\) line \(F_b\mathbb Cv\).\\
 Every \((0,\lambda)\) remains \((0,\lambda)\).};
\end{tikzpicture}
\caption{Exact maps of the original two-test metric and connection,
SC6--9 and SC18--35. The triangle commutes on the full disk for
metrics and on its punctured part for connections. The diagram
is a map diagram, not a spatial plot or a rescaled heat evolution.
The actual time is \(t=t_d+\zeta\), with \(t_d\) given in SC4.
The analytic source is Rodgers--Tao as cited in SC1; the original
test vector and its complete-trace receiver are NC1--2 and WR1--5.}
\label{fig:schur-connection-morphisms}
\end{figure}

\begin{figure}[p]
\centering
\begin{tikzpicture}[>=Stealth,
 box/.style={draw,rounded corners,align=center,text width=5.4cm,
 inner sep=8pt}]
\node[box] (small) at (0,0)
 {Original two-test space \(V\)\\
 \(G=J^TMJ\), \(\operatorname{rank}G_d=1\)\\
 residue \(R\), local reflection \(\mathcal M_b\)};
\node[box] (full) at (7.3,0)
 {Original nineteen-test space \(\mathcal V\)\\
 \(M_d\) invertible, inertia \((18,1,0)\)\\
 \(\Gamma_{\rm full}\) holomorphic, local holonomy \(I_{19}\)};
\draw[->,thick] ([yshift=8pt]small.east)
 --node[above]{\(J\)} ([yshift=8pt]full.west);
\draw[->,thick] ([yshift=-12pt]full.west)
 --node[below]{\(J^\dagger\)} ([yshift=-12pt]small.east);
\node[box,text width=12.7cm] (relation) at (3.65,-3.1)
 {\(\Gamma_{\rm full}J=J\Gamma+W\),\quad
 \(W=(I_{19}-\Pi)\Gamma_{\rm full}J\),\quad
 \(\Pi=JJ^\dagger\)\\[4pt]
 Residues: \(\quad 0=JR+(-JR)\quad\)};
\node[box,text width=12.7cm] at (3.65,-5.8)
 {At \(t_d\):
 \(J(V)\cap J(V)^{\perp_{M_d}}=\mathbb C(Jr)\).\\[4pt]
 \(Jr\) is isotropic, \(M_dJr\ne0\), and
 \(K_\Pi=(Jr)(Jr)^TM_d/s_1\) has
 \(K_\Pi^2=0\), \(K_\Pi\ne0\).};
\end{tikzpicture}
\caption{The exact relation between full and compressed transport,
SC36--48. The two arrows compose to \(I_2\) on the two-test
space and to \(\Pi\) on the full space away from \(t_d\).
The return arrow is meromorphic and its residue is given in
SC44. The complementary term \(W\) is retained; its pole
cancels the compressed pole. The original full theta receiver
is SC1--2 and NC1--4; the certified ambient nondegeneracy is
the complete-theta calculation in
\texttt{UNIFORM\_FULL\_FLAG\_RECEIPT.json}.}
\label{fig:ambient-compressed-cancellation}
\end{figure}
